\begin{theorem}[六倍窗零见证随机命中的精确指数阈值]\label{thm:conclusion-sixfold-zero-random-hit-sharp-threshold}
在定理 \ref{thm:conclusion-sixfold-zero-witness-square-root-law} 的六倍窗设定下，记
\[
p_m:=\frac{|\mathcal Z_m|}{M}.
\]
对任意固定 \(\theta\ge 0\)，定义
\[
\mathsf P_m(\theta):=1-(1-p_m)^{\lfloor M^\theta\rfloor}.
\]
则有严格二相律
\[
\theta<\frac12 \Longrightarrow \mathsf P_m(\theta)\to 0,
\qquad
\theta>\frac12 \Longrightarrow \mathsf P_m(\theta)\to 1.
\]
因此，六倍窗零见证的随机命中阈值在指数尺度上精确等于 \(1/2\)。
\leanverified{paper\_conclusion\_sixfold\_zero\_random\_hit\_sharp\_threshold}
\end{theorem}

\begin{proof}
由定理 \ref{thm:conclusion-sixfold-zero-witness-square-root-law}，
\[
5^{-1/4}M^{-1/2}(1+o(1))\le p_m\le M^{-1/2+o(1)}.
\]
若 \(\theta<1/2\)，则
\[
\lfloor M^\theta\rfloor p_m\le M^{\theta-1/2+o(1)}\to 0.
\]
由并合界，
\[
\mathsf P_m(\theta)\le \lfloor M^\theta\rfloor p_m\to 0.
\]

若 \(\theta>1/2\)，则
\[
\lfloor M^\theta\rfloor p_m\ge 5^{-1/4}M^{\theta-1/2}(1+o(1))\to \infty.
\]
因此
\[
\mathsf P_m(\theta)
=1-(1-p_m)^{\lfloor M^\theta\rfloor}
\ge 1-\exp\!\bigl(-\lfloor M^\theta\rfloor p_m\bigr)\to 1.
\]
证毕。
\end{proof}

\begin{corollary}[固定置信度的最小抽样复杂度为半速率阶]\label{cor:conclusion-sixfold-zero-fixed-confidence-sample-complexity}
\leanverified{paper\_conclusion\_sixfold\_zero\_fixed\_confidence\_sample\_complexity}
对任意固定 \(\eta\in(0,1)\)，定义
\[
N_m(\eta):=\min\{N\ge 1:1-(1-p_m)^N\ge \eta\}.
\]
则
\[
\lim_{\substack{m\to\infty\\ m+2\equiv 0\ (\mathrm{mod}\ 6)}}
\frac{\log N_m(\eta)}{\log M}
=
\frac12.
\]
\end{corollary}

\begin{proof}
任取 \(\varepsilon>0\)。由定理 \ref{thm:conclusion-sixfold-zero-random-hit-sharp-threshold}，
\[
1-(1-p_m)^{M^{1/2-\varepsilon}}\to 0<\eta,
\]
故对充分大的 \(m\) 有
\[
N_m(\eta)>M^{1/2-\varepsilon}.
\]
同理，
\[
1-(1-p_m)^{M^{1/2+\varepsilon}}\to 1>\eta,
\]
故对充分大的 \(m\) 有
\[
N_m(\eta)\le M^{1/2+\varepsilon}.
\]
夹逼并令 \(\varepsilon\downarrow 0\) 即得结论。证毕。
\end{proof}

\begin{proposition}[零见证命中计数的稀疏相与自平均相]\label{prop:conclusion-sixfold-zero-hitcount-sparse-selfaveraging}
设 \(k_1,\dots,k_N\) 为 \(\ZZ/(M\ZZ)\) 上独立均匀样本，并记命中计数
\[
H_m(N):=\sum_{j=1}^{N}\mathbf 1_{\{k_j\in \mathcal Z_m\}}.
\]
则对任意固定 \(\theta\ge 0\)：
\begin{enumerate}
  \item 若 \(\theta<1/2\)，则
  \[
  H_m(M^\theta)\xrightarrow{\mathbb P}0;
  \]
  \item 若 \(\theta>1/2\)，则
  \[
  \frac{H_m(M^\theta)}{\mathbb E\,H_m(M^\theta)}\xrightarrow{\mathbb P}1,
  \qquad
  \mathbb E\,H_m(M^\theta)=M^{\theta-1/2+o(1)}\to\infty.
  \]
\end{enumerate}
因此，指数阈值 \(1/2\) 同时区分了稀疏命中相与自平均命中相。
\end{proposition}

\begin{proof}
随机变量 \(H_m(N)\) 服从二项分布 \(\mathrm{Bin}(N,p_m)\)，故
\[
\mathbb E\,H_m(N)=Np_m,
\qquad
\Var(H_m(N))=Np_m(1-p_m)\le Np_m.
\]
若 \(\theta<1/2\)，则
\[
\mathbb E\,H_m(M^\theta)=M^{\theta-1/2+o(1)}\to 0.
\]
由 Markov 不等式，
\[
\mathbb P\bigl(H_m(M^\theta)\ge 1\bigr)\le \mathbb E\,H_m(M^\theta)\to 0,
\]
故 \(H_m(M^\theta)\xrightarrow{\mathbb P}0\)。

若 \(\theta>1/2\)，则由定理 \ref{thm:conclusion-sixfold-zero-witness-square-root-law}，
\[
\mathbb E\,H_m(M^\theta)\ge 5^{-1/4}M^{\theta-1/2}(1+o(1))\to\infty.
\]
于是对任意 \(\varepsilon>0\)，
\[
\mathbb P\!\left(
\left|H_m(M^\theta)-\mathbb E\,H_m(M^\theta)\right|
>
\varepsilon\,\mathbb E\,H_m(M^\theta)
\right)
\le
\frac{\Var(H_m(M^\theta))}{\varepsilon^2\mathbb E\,H_m(M^\theta)^2}
\le
\frac{1}{\varepsilon^2\mathbb E\,H_m(M^\theta)}
\to 0.
\]
故相对偏差趋于零。证毕。
\end{proof}

\begin{theorem}[同一可见层同时承担两项精确任务时的强逆指数]\label{thm:conclusion-solenoidal-rh-joint-visible-layer-strong-inverse}
固定 \(a>a_{\mathrm c}\)。对每个 \(B\ge 0\)，令 \(\mathcal J_m(B)\) 表示在“至多 \(2^B\) 个可见等价类”的约束下，同一可见证书系统同时承担
\begin{enumerate}
  \item 折叠微态的精确反演；
  \item 冻结相实现点的精确识别
\end{enumerate}
时所能达到的最优统一成功率。更具体地，
\[
\mathcal J_m(B):=
\sup_{|\mathcal C|\le 2^B}
\min\!\bigl\{s_m^{\mathrm{inv}}(\mathcal C),\,s_m^{\mathrm{frz}}(\mathcal C)\bigr\}.
\]
则对任意预算序列 \((B_m)\) 都有
\[
\limsup_{m\to\infty}\frac1m\log_2 \mathcal J_m(B_m)
\le
\limsup_{m\to\infty}\frac{B_m}{m}
-\max\!\bigl\{r_\varphi,g_*^{(2)}\bigr\}.
\]
特别地，若存在 \(\delta>0\) 使
\[
\limsup_{m\to\infty}\frac{B_m}{m}
\le
\max\!\bigl\{r_\varphi,g_*^{(2)}\bigr\}-\delta,
\]
则
\[
\mathcal J_m(B_m)\le 2^{-\delta m+o(m)}.
\]
\end{theorem}

\begin{proof}
对任意可见层 \(\mathcal C\) 且 \(|\mathcal C|\le 2^B\)，都有
\[
\min\!\bigl\{s_m^{\mathrm{inv}}(\mathcal C),s_m^{\mathrm{frz}}(\mathcal C)\bigr\}
\le
s_m^{\mathrm{inv}}(\mathcal C)\le \mathrm{Succ}_m(B),
\]
并且
\[
\min\!\bigl\{s_m^{\mathrm{inv}}(\mathcal C),s_m^{\mathrm{frz}}(\mathcal C)\bigr\}
\le
s_m^{\mathrm{frz}}(\mathcal C)\le P_m^{\mathrm{frz}}(B).
\]
故
\[
\mathcal J_m(B)\le \min\!\bigl\{\mathrm{Succ}_m(B),P_m^{\mathrm{frz}}(B)\bigr\}.
\]
再分别应用定理 \ref{thm:conclusion-solenoidal-rh-fold-strong-converse} 与定理 \ref{thm:conclusion-solenoidal-rh-frozen-visible-strong-converse}，得
\[
\limsup_{m\to\infty}\frac1m\log_2 \mathcal J_m(B_m)
\le
\limsup_{m\to\infty}\frac{B_m}{m}-r_\varphi,
\]
以及
\[
\limsup_{m\to\infty}\frac1m\log_2 \mathcal J_m(B_m)
\le
\limsup_{m\to\infty}\frac{B_m}{m}-g_*^{(2)}.
\]
两式合并即得结论。证毕。
\end{proof}

\begin{theorem}[三预算加法律]\label{thm:conclusion-solenoidal-rh-three-budget-additive-law}
考虑一族 verifier，其在分辨率 \(m\) 上同时使用：
\begin{enumerate}
  \item 一个至多 \(2^{B_m}\) 类的可见证书系统；
  \item \(N_m\) 次对 \(\ZZ/(M\ZZ)\) 的独立均匀随机探测；
  \item 一个有限 prime-stage \(S_m\)。
\end{enumerate}
假设该家族满足：
\begin{enumerate}
  \item 折叠微态反演渐近精确；
  \item 对某个固定 \(a>a_{\mathrm c}\)，冻结相实现点识别渐近精确；
  \item 零见证随机命中概率有非退化下界：
  \[
  \liminf_{m\to\infty}\bigl(1-(1-p_m)^{N_m}\bigr)>0;
  \]
  \item 在定理 \ref{thm:xi-layered-prime-register-sharp-natural-extension} 的意义下，其 hidden prime-ledger 历史含有 \(k(m)\) 层真正分支。
\end{enumerate}
则必有
\[
\liminf_{m\to\infty}
\left(
\frac{B_m}{m}
+
\frac{\log_2 N_m}{m}
\right)
\ge
\max\!\bigl\{r_\varphi,g_*^{(2)}\bigr\}
+
\frac12\log_2\varphi,
\]
并且
\[
|S_m|\ge k(m).
\]
\end{theorem}

\begin{proof}
由定理 \ref{thm:conclusion-solenoidal-rh-double-budget-lower-bound}，前两项可见精确性与 hidden prime-ledger 条件已迫使
\[
\liminf_{m\to\infty}\frac{B_m}{m}\ge \max\!\bigl\{r_\varphi,g_*^{(2)}\bigr\},
\qquad
|S_m|\ge k(m).
\]
另一方面，由随机命中概率的非退化下界，存在常数 \(\eta>0\) 与充分大的 \(m\) 使
\[
1-(1-p_m)^{N_m}\ge \eta.
\]
由推论 \ref{cor:conclusion-sixfold-zero-fixed-confidence-sample-complexity} 可知，这迫使
\[
\frac{\log N_m}{\log M}\ge \frac12-o(1).
\]
又因为
\[
\log_2 M=(\log_2\varphi)\,m+o(m),
\]
故
\[
\liminf_{m\to\infty}\frac{\log_2 N_m}{m}\ge \frac12\log_2\varphi.
\]
将两项下界相加即得所示不等式。prime-stage 下界已如上给出。证毕。
\end{proof}

\begin{corollary}[冻结瓶颈不超过折叠瓶颈时的黄金常数塌缩]\label{cor:conclusion-solenoidal-rh-golden-total-burden-collapse}
若
\[
g_*^{(2)}\le r_\varphi,
\]
则定理 \ref{thm:conclusion-solenoidal-rh-three-budget-additive-law} 简化为
\[
\liminf_{m\to\infty}
\left(
\frac{B_m}{m}
+
\frac{\log_2 N_m}{m}
\right)
\ge
r_\varphi+\frac12\log_2\varphi
=
1-\frac12\log_2\varphi.
\]
特别地，
\[
\liminf_{m\to\infty}
\left(
\frac{B_m}{m}
+
\frac{\log_2 N_m}{m}
\right)
\ge
0.6528790431846913\ldots
\]
\end{corollary}

\begin{proof}
在定理 \ref{thm:conclusion-solenoidal-rh-three-budget-additive-law} 中代入
\[
\max\!\bigl\{r_\varphi,g_*^{(2)}\bigr\}=r_\varphi,
\]
再用
\[
r_\varphi=\log_2\frac{2}{\varphi}=1-\log_2\varphi
\]
即可。证毕。
\end{proof}

\begin{theorem}[zero-count 证书的碰撞区间具有单位加法宽度与整一幂乘法宽度]\label{thm:conclusion-zero-count-certificate-unit-additive-width}
在六倍窗设定下，令
\[
L_m:=\frac{1+2C_\varphi^2+o(1)}{M},
\qquad
U_m:=1-M^{-1/2+o(1)}.
\]
其中 \(L_m\) 是真实碰撞的黄金共振下界，而 \(U_m\) 是由 zero-count 证书导出的碰撞上界。则区间
\[
I_m^{\mathrm{zc}}:=[L_m,U_m]
\]
满足
\[
U_m-L_m=1-o(1),
\]
以及
\[
\frac{U_m}{L_m}
=
\frac{M^{1+o(1)}}{1+2C_\varphi^2}.
\]
特别地，
\[
\lim_{\substack{m\to\infty\\ m+2\equiv0\ (\mathrm{mod}\ 6)}}
\frac{\log(U_m/L_m)}{\log M}=1.
\]
因此，即使完整给定 \(|\mathcal Z_m|\)，所导出的碰撞区间在加法尺度上几乎与 \([0,1]\) 同宽，而在乘法尺度上仍比真实 \(M^{-1}\) 层级粗糙整整一个 \(M^{1+o(1)}\) 幂阶。
\end{theorem}

\begin{proof}
由定理 \ref{thm:conclusion-zero-count-collision-blindness}，
\[
U_m=1-M^{-1/2+o(1)},
\qquad
L_m=\frac{1+2C_\varphi^2+o(1)}{M}.
\]
于是
\[
U_m-L_m
=
1-M^{-1/2+o(1)}-\frac{1+2C_\varphi^2+o(1)}{M}
=
1-o(1).
\]
又因为 \(U_m\to 1\)，故
\[
\frac{U_m}{L_m}
=
\frac{1-M^{-1/2+o(1)}}{(1+2C_\varphi^2+o(1))/M}
=
\frac{M^{1+o(1)}}{1+2C_\varphi^2}.
\]
两边取对数并除以 \(\log M\)，即可得到指数极限。证毕。
\end{proof}

\begin{conjecture}[三坐标预算正交性]\label{conj:conclusion-solenoidal-rh-three-coordinate-budget-orthogonality}
在任何可量化正规化的 solenoidal 审计框架内，若一族 verifier 同时要求：
\begin{enumerate}
  \item 折叠微态反演渐近精确；
  \item 冻结相实现点识别渐近精确；
  \item 六倍窗零见证随机命中概率有正下界；
  \item hidden prime-ledger 历史含有 \(k(m)\) 层真正分支，
\end{enumerate}
则其资源三元组
\[
\left(
\frac{B_m}{m},
\frac{\log_2 N_m}{m},
|S_m|
\right)
\]
在渐近上满足不可交换的坐标下界
\[
\frac{B_m}{m}\ge \max\!\bigl\{r_\varphi,g_*^{(2)}\bigr\}-o(1),
\qquad
\frac{\log_2 N_m}{m}\ge \frac12\log_2\varphi-o(1),
\qquad
|S_m|\ge k(m),
\]
并且不存在任何非平凡 trade-off 原理，可用一坐标的超额去补偿另一坐标的亏缺。
等价地，本文框架中的 no-go 区域并非由单一标量预算切出，而是一个真正的正交角域：visible witness rate、square-root sampling exponent 与 cofinal prime-stage depth 分别构成三条彼此不可替代的临界轴。
\end{conjecture}
